\documentclass[aspectratio=169,utf8,t]{beamer}

% Idioma e Codificacao
\usepackage[brazilian]{babel}
\usepackage{booktabs}
\usepackage{array}
\usepackage{listings}
\usepackage[normalem]{ulem}
\usepackage{tikz}
\usetikzlibrary{shapes,arrows,positioning,calc,fit}

% Carrega o Tema Modelo dos Scouts (Widescreen 16:9, Nunito Sans)
\usetheme{scouts}

% Configuracao para Codigo-Fonte (Python)
\lstset{
    basicstyle=\ttfamily\scriptsize,
    keywordstyle=\bfseries\color{ScoutNavy},
    stringstyle=\color{ScoutGreen},
    commentstyle=\itshape\color{gray},
    backgroundcolor=\color{black!3},
    frame=single,
    rulecolor=\color{black!15},
    breaklines=true,
    showstringspaces=false,
    tabsize=2
}

% Metadados Institucionais
\author{Prof. Me. Reginaldo Fernandes}
\institute{Instituto Federal do Ceará -- Campus Tauá \\ Tecnologia em Telemática}
\title{Banco de Dados}
\subtitle{Aula 06: Entidades Fracas e Modelo EER (Especialização/Generalização)}
\date{14 de Outubro de 2026}

\begin{document}

% ------------------------------------------------------------------------------
% Capa da Apresentacao
% ------------------------------------------------------------------------------
\changecolours[inverse]{ScoutGreen}
\frame[plain]{\titlepage}

% ------------------------------------------------------------------------------
% Sumario Executivo da Aula
% ------------------------------------------------------------------------------
\changecolours{ScoutGreen}
\begin{frame}{Agenda do Encontro}
  \begin{columns}[t]
    \begin{column}{0.48\textwidth}
      \scriptsize
      \tableofcontents[sections={1-5}]
    \end{column}
    \begin{column}{0.48\textwidth}
      \scriptsize
      \tableofcontents[sections={6-9}]
    \end{column}
  \end{columns}
\end{frame}

% ------------------------------------------------------------------------------
% Modulos Tematicos da Aula
% ------------------------------------------------------------------------------
\changecolours{ScoutGreen}
\section[Revisão e Objetivos]{Ponte Pedagógica e Objetivos de Aprendizagem}

% ------------------------------------------------------------------------------
\twocolframe[title={Retrospectiva da Aula 05 e o Novo Desafio},leftcol=7.1cm,rightcol=6.9cm]{%
  \alert{O Que Consolidamos na Aula 05?}\par
  \itemseps[topsepi=0.5mm,itemsepi=1mm,labelsep=2mm,leftmargini=3mm]
  \begin{itemize}
    \item \textbf{Relacionamentos no MER:} Associações semânticas como subconjunto do produto cartesiano ($R \subseteq E_1 \times E_2$).
    \item \textbf{Cardinalidade:} Razões $1:1$, $1:N$ e $N:M$ derivadas estritamente das regras de negócio.
    \item \textbf{Participação:} Linha dupla (total/obrigatória) e simples (parcial/opcional).
    \item \textbf{Atributos Descritivos:} Alocados no losango de ligação.
  \end{itemize}
}{%
  \alert{O Desafio Central da Aula 06}\par
  \itemseps[topsepi=0.5mm,itemsepi=1mm,labelsep=2mm,leftmargini=3mm]
  \begin{itemize}
    \item \textbf{Entidades Dependentes:} E quando uma entidade do mundo real \textbf{não possui atributos suficientes} para formar uma chave primária própria e autônoma?
      \begin{itemize}\setlength{\itemsep}{0.3mm}
        \item Ex: As portas de rede de um roteador.
        \item Ex: Os dependentes de um funcionário.
      \end{itemize}
    \item \textbf{Taxonomias e Herança:} Como representar hierarquias do tipo ``é um'' sem duplicar atributos comuns?
  \end{itemize}
}

% ------------------------------------------------------------------------------
\twocolframe[title={Objetivos de Aprendizagem da Aula 06},leftcol=7.1cm,rightcol=6.9cm]{%
  \alert{Competências Teóricas e Conceituais}\par
  \itemseps[topsepi=0.5mm,itemsepi=1mm,labelsep=2mm,leftmargini=3mm]
  \begin{itemize}
    \item \textbf{Entidades Fracas:} Dominar os conceitos de dependência existencial e de identificação.
    \item \textbf{Chave Parcial (Discriminador):} Formar a chave primária composta com a entidade proprietária.
    \item \textbf{Notação de Peter Chen:} Aplicar retângulos duplos, losangos duplos e sublinhado tracejado.
    \item \textbf{Modelo EER:} Projetar hierarquias com superclasses, subclasses e herança de atributos.
  \end{itemize}
}{%
  \alert{Competências Metodológicas e Práticas}\par
  \itemseps[topsepi=0.5mm,itemsepi=1mm,labelsep=2mm,leftmargini=3mm]
  \begin{itemize}
    \item \textbf{Restrições Estruturais EER:} Classificar especializações em Disjuntas ($d$) ou Sobrepostas ($o$), Totais ou Parciais.
    \item \textbf{Modelagem em Telemática:} Aplicar conceitos a equipamentos de rede, interfaces e topologias.
    \item \textbf{Laboratório Computacional:} Validar regras em memória via metamodelo em Python orientado a objetos.
    \item \textbf{Alinhamento AVP1:} Esboçar entidades fracas e subclasses para o projeto semestral.
  \end{itemize}
}

\changecolours{ScoutGreen}
\section[Entidades Fracas]{O Conceito Formal de Entidade Fraca no MER}

% ------------------------------------------------------------------------------
\twocolframe[title={Definição Formal de Entidade Fraca (Peter Chen, 1976)},leftcol=7.1cm,rightcol=6.9cm]{%
  \alert{O Que é uma Entidade Fraca?}\par
  \itemseps[topsepi=0.5mm,itemsepi=1mm,labelsep=2mm,leftmargini=3mm]
  \begin{itemize}
    \item \textbf{Definição Clássica:} É uma entidade que \textbf{não possui atributos próprios suficientes} para compor uma chave primária autônoma e exclusiva no sistema.
    \item \textbf{Entidade Forte / Regular (*Owner*):}
      Possui chave primária intrínseca e existência independente no minimundo.
    \item \textbf{Subordinação Semântica:}
      A entidade fraca só adquire identidade e sentido quando contextualizada por meio de uma entidade proprietária.
  \end{itemize}
}{%
  \alert{Tipos de Dependência Estrutural}\par
  \itemseps[topsepi=0.5mm,itemsepi=1mm,labelsep=2mm,leftmargini=3mm]
  \begin{itemize}
    \item \textbf{Dependência Existencial:}
      Uma instância da entidade fraca não pode existir no banco se a entidade proprietária correspondente for excluída.
    \item \textbf{Dependência de Identificação:}
      A chave primária da entidade fraca é obrigatoriamente formada pela combinação da chave da proprietária com seu próprio discriminador local.
    \item \textbf{Regra Canônica:}
      Toda dependência de identificação implica dependência existencial, mas o inverso nem sempre é verdadeiro!
  \end{itemize}
}

% ------------------------------------------------------------------------------
\begin{frame}{Gramática Gráfica Canônica de Peter Chen para Entidades Fracas}
  \vspace{0.1cm}
  \begin{center}
    \tikz[scale=0.76, transform shape]{
      % Entidade Forte (Proprietaria)
      \node[draw=ScoutGreen, fill=ScoutGreen!15, thick, font=\bfseries\sffamily\scriptsize, minimum width=2.4cm, minimum height=0.7cm] (forte) {FUNCIONARIO};
      \node[draw=ScoutNavy, fill=ScoutNavy!10, ellipse, font=\tiny\sffamily, above left=0.5cm and 0.1cm of forte] (pk_forte) {\underline{cpf}};
      \draw[thick, ScoutNavy] (forte) -- (pk_forte);

      % Relacionamento Identificador (Losango Duplo)
      \node[diamond, draw=ScoutNavy, double, double distance=1.2pt, fill=ScoutNavy!15, thick, inner sep=2pt, font=\bfseries\sffamily\tiny, right=1.2cm of forte] (rel_id) {POSSUI\_DEP};

      % Entidade Fraca (Retangulo Duplo)
      \node[draw=ScoutGreen, double, double distance=1.2pt, fill=ScoutGreen!15, thick, font=\bfseries\sffamily\scriptsize, minimum width=2.4cm, minimum height=0.7cm, right=1.2cm of rel_id] (fraca) {DEPENDENTE};
      \node[draw=ScoutNavy, fill=ScoutNavy!10, ellipse, font=\tiny\sffamily, above right=0.5cm and 0.1cm of fraca] (disc) {\dashuline{nome\_dep}};
      \draw[thick, ScoutNavy] (fraca) -- (disc);

      % Linhas de conexao
      \draw[thick, black!80] (forte) -- node[above, font=\scriptsize\bfseries] {1} (rel_id);
      \draw[thick, double, double distance=1.5pt, black!80] (rel_id) -- node[above, font=\scriptsize\bfseries] {N} (fraca);
    }
  \end{center}
  \vspace{0.15cm}
  \begin{columns}[t]
    \begin{column}{0.48\textwidth}
      \footnotesize
      \alert{Convenções Visuais Obrigatórias:}\par
      $\bullet$ \textbf{Retângulo Duplo:} Indica a entidade fraca.\\
      $\bullet$ \textbf{Losango Duplo:} Indica o relacionamento identificador.\\
      $\bullet$ \textbf{Linha Dupla Fraca-Losango:} Participação total obrigatória.
    \end{column}
    \begin{column}{0.48\textwidth}
      \footnotesize
      \alert{Semântica da Ligação:}\par
      $\bullet$ O relacionamento identificador é sempre $1:N$ do lado do proprietário para o fraco.\\
      $\bullet$ O discriminador (\texttt{nome\_dep}) é grafado com sublinhado tracejado.
    \end{column}
  \end{columns}
\end{frame}

\changecolours{ScoutGreen}
\section[Chave Parcial]{Chave Parcial (Discriminador) e Chaves Compostas}

% ------------------------------------------------------------------------------
\twocolframe[title={O Conceito de Chave Parcial (Discriminador)},leftcol=7.1cm,rightcol=6.9cm]{%
  \alert{O Que é a Chave Parcial?}\par
  \itemseps[topsepi=0.5mm,itemsepi=1mm,labelsep=2mm,leftmargini=3mm]
  \begin{itemize}
    \item \textbf{Definição:} Conjunto de atributos que distingue diferentes instâncias da entidade fraca \textbf{que estão associadas a uma mesma entidade proprietária}.
    \item \textbf{Unicidade Local vs. Global:}
      \begin{itemize}\setlength{\itemsep}{0.3mm}
        \item Localmente: É única apenas para aquele proprietário.
        \item Globalmente: Pode se repetir livremente sob proprietários distintos.
      \end{itemize}
    \item \textbf{Notação Gráfica:} Elipse conectada à entidade fraca com o nome do atributo \textbf{sublinhado por linha tracejada}.
  \end{itemize}
}{%
  \alert{Formação da Chave Composta Completa}\par
  \itemseps[topsepi=0.5mm,itemsepi=1mm,labelsep=2mm,leftmargini=3mm]
  \begin{itemize}
    \item \textbf{Equação da Identificação:}
      $$\text{PK}(\text{Fraca}) = \text{PK}(\text{Proprietária}) \cup \text{Discriminador}$$
    \item \textbf{Exemplo Corporativo:}
      \begin{itemize}\setlength{\itemsep}{0.3mm}
        \item Proprietária: \texttt{FUNCIONARIO}(\underline{\texttt{cpf}})
        \item Fraca: \texttt{DEPENDENTE}(\dashuline{\texttt{nome\_dep}})
        \item Chave Lógica: (\texttt{cpf}, \texttt{nome\_dep})
      \end{itemize}
    \item \textbf{Integridade Referencial:}
      No modelo relacional futuro, a PK da proprietária migra como Foreign Key (FK) compondo a Primary Key (PK).
  \end{itemize}
}

% ------------------------------------------------------------------------------
\begin{frame}{Exemplo Cotidiano: Redes de Hotéis e Quartos}
  \vspace{0.1cm}
  \begin{center}
    \tikz[scale=0.70, transform shape]{
      % Hotel (Entidade Forte)
      \node[draw=ScoutGreen, fill=ScoutGreen!15, thick, font=\bfseries\sffamily\scriptsize, minimum width=2.4cm, minimum height=0.65cm] (hotel) {HOTEL};
      \node[draw=ScoutNavy, fill=ScoutNavy!10, ellipse, font=\tiny\sffamily, above left=0.5cm and 0.1cm of hotel] (pk_hotel) {\underline{id\_hotel}};
      \node[draw=ScoutNavy, fill=ScoutNavy!10, ellipse, font=\tiny\sffamily, below left=0.4cm and 0.1cm of hotel] (nome_h) {nome\_hotel};
      \draw[thick, ScoutNavy] (hotel) -- (pk_hotel);
      \draw[thick, ScoutNavy] (hotel) -- (nome_h);

      % Relacionamento Identificador
      \node[diamond, draw=ScoutNavy, double, double distance=1.2pt, fill=ScoutNavy!15, thick, inner sep=2pt, font=\bfseries\sffamily\tiny, right=1.0cm of hotel] (rel) {POSSUI\_QUARTO};

      % Quarto (Entidade Fraca)
      \node[draw=ScoutGreen, double, double distance=1.2pt, fill=ScoutGreen!15, thick, font=\bfseries\sffamily\scriptsize, minimum width=2.4cm, minimum height=0.65cm, right=1.0cm of rel] (quarto) {QUARTO};
      \node[draw=ScoutNavy, fill=ScoutNavy!10, ellipse, font=\tiny\sffamily, above right=0.5cm and 0.1cm of quarto] (disc) {\dashuline{numero\_quarto}};
      \node[draw=ScoutNavy, fill=ScoutNavy!10, ellipse, font=\tiny\sffamily, below right=0.4cm and 0.1cm of quarto] (valor) {valor\_diaria};
      \draw[thick, ScoutNavy] (quarto) -- (disc);
      \draw[thick, ScoutNavy] (quarto) -- (valor);

      % Conexoes Chen
      \draw[thick, black!80] (hotel) -- node[above, font=\scriptsize\bfseries] {1} (rel);
      \draw[thick, double, double distance=1.5pt, black!80] (rel) -- node[above, font=\scriptsize\bfseries] {N} (quarto);
    }
  \end{center}
  \vspace{-0.1cm}
  \begin{table}[h]
    \centering
    \resizebox{0.78\textwidth}{!}{%
    \begin{tabular}{l l l l}
      \toprule
      \textbf{ID Hotel (Proprietário)} & \textbf{Nome do Hotel} & \textbf{Discriminador (Quarto)} & \textbf{Chave Composta Global} \\
      \midrule
      \texttt{H-01} & \texttt{Hotel Beira-Mar} & \texttt{101} & \texttt{(H-01, 101)} \\
      \texttt{H-01} & \texttt{Hotel Beira-Mar} & \texttt{102} & \texttt{(H-01, 102)} \\
      \texttt{H-02} & \texttt{Pousada do Sol} & \texttt{101} & \texttt{(H-02, 101)} \\
      \bottomrule
    \end{tabular}%
    }
  \end{table}
  \vspace{-0.1cm}
  {\centering\tiny \textbf{Observação:} O quarto 101 repete-se em hotéis distintos sem qualquer colisão de chave primária.\par}
\end{frame}

\changecolours{ScoutGreen}
\section[Modelo EER]{O Modelo Entidade-Relacionamento Estendido (EER)}

% ------------------------------------------------------------------------------
\twocolframe[title={Por Que Estender o MER Clássico?},leftcol=7.1cm,rightcol=6.9cm]{%
  \alert{As Limitações do MER Tradicional}\par
  \itemseps[topsepi=0.5mm,itemsepi=1mm,labelsep=2mm,leftmargini=3mm]
  \begin{itemize}
    \item \textbf{Complexidade Crescente:} Sistemas modernos exigem modelar relações semânticas do tipo ``é um'' (\textit{is-a}).
    \item \textbf{Problema da Redundância:} Sem herança, atributos comuns (como \texttt{nome}, \texttt{fabricante}, \texttt{data\_compra}) seriam replicados em dezenas de tabelas.
    \item \textbf{O Modelo EER:} Introduzido para agregar conceitos de orientação a objetos, abstração e classificação taxonômica.
  \end{itemize}
}{%
  \alert{Superclasses e Subclasses}\par
  \itemseps[topsepi=0.5mm,itemsepi=1mm,labelsep=2mm,leftmargini=3mm]
  \begin{itemize}
    \item \textbf{Superclasse:} Tipo de entidade genérico que agrupa características comuns a várias entidades distintas.
    \item \textbf{Subclasse:} Agrupamento especializado de instâncias que compartilham atributos ou relacionamentos próprios exclusivos.
    \item \textbf{Exemplo Intuitivo:}
      \begin{itemize}\setlength{\itemsep}{0.3mm}
        \item Superclasse: \texttt{DISPOSITIVO\_REDE}
        \item Subclasses: \texttt{ROTEADOR}, \texttt{SWITCH}, \texttt{FIREWALL}
      \end{itemize}
  \end{itemize}
}

% ------------------------------------------------------------------------------
\twocolframe[title={O Mecanismo de Herança de Tipos no EER},leftcol=7.1cm,rightcol=6.9cm]{%
  \alert{Herança Universal de Atributos}\par
  \itemseps[topsepi=0.5mm,itemsepi=1mm,labelsep=2mm,leftmargini=3mm]
  \begin{itemize}
    \item \textbf{Regra Áurea:} Toda instância que pertence a uma subclasse herda \textbf{automaticamente todos os atributos} da sua superclasse.
    \item \textbf{Atributos Específicos:}
      Além dos atributos herdados, a subclasse pode possuir atributos descritivos que só fazem sentido para ela.
    \item \textbf{Identidade Única:}
      A chave primária da subclasse é sempre a mesma chave primária herdada da superclasse!
  \end{itemize}
}{%
  \alert{Herança de Relacionamentos}\par
  \itemseps[topsepi=0.5mm,itemsepi=1mm,labelsep=2mm,leftmargini=3mm]
  \begin{itemize}
    \item \textbf{Participação Indireta:}
      A subclasse participa automaticamente de qualquer relacionamento no qual sua superclasse esteja envolvida.
    \item \textbf{Relacionamentos Exclusivos:}
      A subclasse pode também participar de relacionamentos especializados que não se aplicam às demais irmãs.
    \item \textbf{Exemplo:}
      Todo \texttt{ROTEADOR} é um \texttt{DISPOSITIVO} (está em um rack), mas só roteadores participam de \texttt{TABELA\_BGP}.
  \end{itemize}
}

\changecolours{ScoutGreen}
\section[Especialização e Generalização]{Especialização, Generalização e a Notação EER}

% ------------------------------------------------------------------------------
\twocolframe[title={Abordagens: Top-Down vs. Bottom-Up},leftcol=7.1cm,rightcol=6.9cm]{%
  \alert{Especialização (Abordagem Top-Down)}\par
  \itemseps[topsepi=0.5mm,itemsepi=1mm,labelsep=2mm,leftmargini=3mm]
  \begin{itemize}
    \item \textbf{Processo:} Parte-se de um tipo de entidade genérico existente e criam-se subclasses baseadas em características distintivas.
    \item \textbf{Quando Ocorre?}
      Quando o analista percebe que certos subgrupos de instâncias possuem propriedades específicas ou regras de negócio exclusivas.
    \item \textbf{Exemplo:} \texttt{FUNCIONARIO} especializa-se em \texttt{ENGENHEIRO}, \texttt{SECRETARIO} e \texttt{TECNICO}.
  \end{itemize}
}{%
  \alert{Generalização (Abordagem Bottom-Up)}\par
  \itemseps[topsepi=0.5mm,itemsepi=1mm,labelsep=2mm,leftmargini=3mm]
  \begin{itemize}
    \item \textbf{Processo:} Identificam-se múltiplos tipos de entidades já existentes com características comuns e cria-se uma superclasse genérica para unificá-los.
    \item \textbf{Quando Ocorre?}
      Em fases de consolidação de requisitos onde modelos parciais são integrados.
    \item \textbf{Exemplo:} Unificar \texttt{CARRO}, \texttt{CAMINHAO} e \texttt{MOTO} na superclasse \texttt{VEICULO}.
  \end{itemize}
}

% ------------------------------------------------------------------------------
\begin{frame}{Representação Gráfica Canônica no Diagrama EER}
  \vspace{0.1cm}
  \begin{center}
    \tikz[scale=0.75, transform shape]{
      % Superclasse
      \node[draw=ScoutGreen, fill=ScoutGreen!15, thick, font=\bfseries\sffamily\scriptsize, minimum width=3.0cm, minimum height=0.65cm] (super) {VEICULO};
      \node[draw=ScoutNavy, fill=ScoutNavy!10, ellipse, font=\tiny\sffamily, above left=0.4cm and 0.3cm of super] (pk) {\underline{placa}};
      \node[draw=ScoutNavy, fill=ScoutNavy!10, ellipse, font=\tiny\sffamily, above right=0.4cm and 0.3cm of super] (fab) {marca};
      \draw[thick, ScoutNavy] (super) -- (pk);
      \draw[thick, ScoutNavy] (super) -- (fab);

      % Circulo EER (Letra d indicando disjuncao)
      \node[draw=ScoutNavy, fill=ScoutNavy!20, circle, thick, inner sep=2.5pt, font=\bfseries\sffamily\scriptsize, below=0.65cm of super] (circ) {d};

      % Subclasses
      \node[draw=ScoutGreen, fill=ScoutGreen!15, thick, font=\bfseries\sffamily\scriptsize, minimum width=2.2cm, minimum height=0.65cm, below left=0.7cm and 1.0cm of circ] (sub1) {CARRO};
      \node[draw=ScoutNavy, fill=ScoutNavy!10, ellipse, font=\tiny\sffamily, below=0.35cm of sub1] (atr_sub1) {qtd\_portas};
      \draw[thick, ScoutNavy] (sub1) -- (atr_sub1);

      \node[draw=ScoutGreen, fill=ScoutGreen!15, thick, font=\bfseries\sffamily\scriptsize, minimum width=2.2cm, minimum height=0.65cm, below right=0.7cm and 1.0cm of circ] (sub2) {CAMINHAO};
      \node[draw=ScoutNavy, fill=ScoutNavy!10, ellipse, font=\tiny\sffamily, below=0.35cm of sub2] (atr_sub2) {capacidade\_carga};
      \draw[thick, ScoutNavy] (sub2) -- (atr_sub2);

      % Linhas conectando superclasse ao circulo e circulo as subclasses
      \draw[thick, double, double distance=1.5pt, black!80] (super) -- (circ);
      \draw[thick, black!80] (circ) -| (sub1);
      \draw[thick, black!80] (circ) -| (sub2);

      % Simbolo de subconjunto U na ligacao
      \node[font=\footnotesize, left=0.2cm of sub1.north, yshift=0.3cm] {$\cup$};
      \node[font=\footnotesize, right=0.2cm of sub2.north, yshift=0.3cm] {$\cup$};
    }
  \end{center}
  \vspace{0.05cm}
  \centering{\tiny \textbf{Notação EER:} A linha dupla para o círculo indica especialização total; a letra ``d'' indica disjunção exclusiva.}
\end{frame}

\changecolours{ScoutGreen}
\section[Restrições Estruturais]{Restrições Estruturais: Disjunção e Integralidade}

% ------------------------------------------------------------------------------
\twocolframe[title={Restrições de Disjunção: Disjunta (d) vs. Sobreposta (o)},leftcol=7.1cm,rightcol=6.9cm]{%
  \alert{Especialização Disjunta (Letra $d$)}\par
  \itemseps[topsepi=0.5mm,itemsepi=1mm,labelsep=2mm,leftmargini=3mm]
  \begin{itemize}
    \item \textbf{Regra de Exclusividade:} Uma entidade da superclasse pode pertencer a \textbf{no máximo uma} subclasse especializada.
    \item \textbf{Subconjuntos Disjuntos:}
      As subclasses são mutuamente exclusivas ($S_1 \cap S_2 = \emptyset$).
    \item \textbf{Exemplo:} Um equipamento em operação na rede atua como \texttt{ROTEADOR} OU como \texttt{SWITCH}, mas nunca como ambos simultaneamente.
  \end{itemize}
}{%
  \alert{Especialização Sobreposta (Letra $o$)}\par
  \itemseps[topsepi=0.5mm,itemsepi=1mm,labelsep=2mm,leftmargini=3mm]
  \begin{itemize}
    \item \textbf{Regra de Não-Exclusividade:} Uma mesma entidade pode pertencer a \textbf{múltiplas subclasses} ao mesmo tempo.
    \item \textbf{Interseção Permitida:}
      $S_1 \cap S_2 \neq \emptyset$.
    \item \textbf{Exemplo Acadêmico:} Em uma universidade, uma mesma \texttt{PESSOA} pode ser cadastrada simultaneamente como \texttt{ALUNO} e como \texttt{FUNCIONARIO}.
  \end{itemize}
}

% ------------------------------------------------------------------------------
\twocolframe[title={Restrições de Integralidade e Matriz Canônica EER},leftcol=7.1cm,rightcol=6.9cm]{%
  \alert{Integralidade (Completude) da Hierarquia}\par
  \itemseps[topsepi=0.5mm,itemsepi=1mm,labelsep=2mm,leftmargini=3mm]
  \begin{itemize}
    \item \textbf{Total (Linha Dupla no EER):}
      Toda instância da superclasse \textbf{deve obrigatoriamente} pertencer a pelo menos uma das subclasses. Não existe superclasse ``pura''.
    \item \textbf{Parcial (Linha Simples no EER):}
      Uma instância da superclasse pode existir sem ser classificada em nenhuma das subclasses especializadas.
  \end{itemize}
}{%
  \alert{As Quatro Combinações Canônicas}\par
  \vspace{0.1cm}
  \begin{table}[h]
    \centering
    \resizebox{\linewidth}{!}{%
    \begin{tabular}{l l l}
      \toprule
      \textbf{Integralidade} & \textbf{Disjunta ($d$)} & \textbf{Sobreposta ($o$)} \\
      \midrule
      \textbf{Total (Dupla)} & Exclusiva e Obrigatória & Múltipla e Obrigatória \\
      \addlinespace[1.5mm]
      \textbf{Parcial (Simples)} & Exclusiva e Opcional & Múltipla e Opcional \\
      \bottomrule
    \end{tabular}%
    }
  \end{table}
  \vspace{0.15cm}
  \begin{itemize}\setlength{\itemsep}{0.5mm}
    \item \textbf{Importância de Projeto:}
      Determina como as tabelas serão mapeadas no modelo lógico relacional (tabela única, tabelas separadas ou tabela base).
  \end{itemize}
}

\changecolours{ScoutGreen}
\section[Laboratório em Python]{Laboratório Computacional: Metamodelo EER em Python}

% ------------------------------------------------------------------------------
\begin{frame}[fragile]{Código Python: Estrutura da Entidade Fraca e Discriminador}
  \vspace{0.1cm}
  \begin{columns}[t]
    \begin{column}{0.48\textwidth}
      \footnotesize
      \alert{Entidade Fraca e Chave Composta}\par
      \vspace{0.15cm}
      \begin{itemize}\setlength{\itemsep}{1.2mm}\setlength{\parskip}{0pt}
        \item A entidade forte possui chave primaria autonoma (\texttt{id\_hotel}).
        \item A entidade fraca armazena apenas seu discriminador (\texttt{numero\_quarto}).
        \item A propriedade \texttt{chave\_completa} calcula dinamicamente a tupla identificadora global.
      \end{itemize}
    \end{column}

    \begin{column}{0.48\textwidth}
      \alert{Implementação Orientada a Objetos}\par
      \vspace{0.1cm}
\begin{lstlisting}[language=Python,basicstyle=\ttfamily\fontsize{5.2}{6.2}\selectfont]
class Quarto:
    def __init__(self, numero_quarto, tipo_leito, diaria):
        # Discriminador (Chave Parcial)
        self.numero_quarto = numero_quarto
        self.tipo_leito = tipo_leito
        self.diaria = diaria
        self.hotel = None

    @property
    def chave_completa(self):
        if not self.hotel:
            return None
        # PK Composta: (PK_Proprietario, Discriminador)
        return (self.hotel.id_hotel, 
                self.numero_quarto)
\end{lstlisting}
    \end{column}
  \end{columns}
\end{frame}

% ------------------------------------------------------------------------------
\begin{frame}[fragile]{Código Python: Herança e Bloqueio de Disjunção no EER}
  \vspace{0.1cm}
  \begin{columns}[t]
    \begin{column}{0.48\textwidth}
      \footnotesize
      \alert{Herança e Validação Disjunta}\par
      \vspace{0.15cm}
      \begin{itemize}\setlength{\itemsep}{1.2mm}\setlength{\parskip}{0pt}
        \item A superclasse gerencia atributos universais (\texttt{placa}, \texttt{marca}, \texttt{modelo}).
        \item Subclasses herdam via \texttt{super().\_\_init\_\_}.
        \item O validador rejeita cadastros simultaneos quando a regra de disjuncao for estrita ($d$).
      \end{itemize}
    \end{column}

    \begin{column}{0.48\textwidth}
      \alert{Herança e Trava Semântica}\par
      \vspace{0.1cm}
\begin{lstlisting}[language=Python,basicstyle=\ttfamily\fontsize{5.0}{6.0}\selectfont]
class Carro(Veiculo):
    def __init__(self, placa, marca, mod, portas, mala):
        super().__init__(placa, marca, mod)
        self.qtd_portas = portas
        self.porta_malas_litros = mala

def especializar(validador, placa, subclasse):
    subs = validador.classificacoes[placa]
    if validador.disjuncao == 'd' and len(subs) > 0:
        raise ValueError(f"Violacao (d): {placa} ja e {list(subs)[0]}!")
    subs.add(subclasse)
\end{lstlisting}
    \end{column}
  \end{columns}
\end{frame}

\changecolours{ScoutGreen}
\section[Conclusão e Próximos Passos]{Conclusão, Atividades e Conexão com a Aula 07}

% ------------------------------------------------------------------------------
\twocolframe[title={Síntese Conceitual da Aula 06},leftcol=7.1cm,rightcol=6.9cm]{%
  \alert{Checklist de Competências Consolidadas}\par
  \itemseps[topsepi=0.5mm,itemsepi=1mm,labelsep=2mm,leftmargini=3mm]
  \begin{itemize}
    \item \textbf{Entidade Fraca:} Ausência de chave própria e dependência existencial/identificação.
    \item \textbf{Discriminador:} Chave parcial grafada com sublinhado tracejado (\dashuline{atributo}).
    \item \textbf{Gramática Chen:} Retângulo duplo, losango duplo e linha dupla para a entidade fraca.
    \item \textbf{Modelo EER:} Superclasses, subclasses e herança universal de atributos e relacionamentos.
    \item \textbf{Restrições EER:} Disjunção ($d/o$) e Integralidade (Total/Parcial).
  \end{itemize}
}{%
  \alert{Boas Práticas de Modelagem Conceitual}\par
  \itemseps[topsepi=0.5mm,itemsepi=1mm,labelsep=2mm,leftmargini=3mm]
  \begin{itemize}
    \item \textbf{Não Abusar de Entidades Fracas:}
      Só utilize entidade fraca se a existência e identificação forem estritamente dependentes de outra entidade do minimundo.
    \item \textbf{Clareza na Especialização:}
      Não crie subclasses se não houver atributos específicos ou relacionamentos exclusivos justificando a divisão.
    \item \textbf{Documentação Formal:}
      Sempre explicite se a hierarquia é ($d$) ou ($o$), e se a participação na superclasse é total ou parcial.
  \end{itemize}
}

% ------------------------------------------------------------------------------
\twocolframe[title={Ponte Pedagógica: Aula 07 (Sábado Letivo)},leftcol=7.1cm,rightcol=6.9cm]{%
  \alert{O Próximo Encontro: Aula 07 (17/10/2026)}\par
  \itemseps[topsepi=0.5mm,itemsepi=1mm,labelsep=2mm,leftmargini=3mm]
  \begin{itemize}
    \item \textbf{Sábado Letivo -- Estudo Dirigido AVP1:}
      Desenvolvimento do projeto prático de modelagem conceitual de banco de dados.
    \item \textbf{Escopo do Projeto Semestral:}
      Definição do minimundo de aplicação (gestão de frotas, reservas hoteleiras ou controle comercial).
    \item \textbf{Modelagem Hands-on:}
      Construção completa do DER conceitual aplicando entidades regulares, fracas, atributos e relacionamentos.
  \end{itemize}
}{%
  \alert{Atividades de Fixação da Aula 06}\par
  \itemseps[topsepi=0.5mm,itemsepi=1mm,labelsep=2mm,leftmargini=3mm]
  \begin{itemize}
    \item \textbf{Exercícios Conceituais:}
      Resolver as 5 questões analíticas do plano de aula no caderno de estudos.
    \item \textbf{Jupyter Notebook Interativo:}
      Executar o script Python para testar os bloqueios de discriminador duplicado e disjunção.
    \item \textbf{Esboço da AVP1:}
      Identificar ao menos uma entidade fraca ou hierarquia EER no escopo da sua equipe.
  \end{itemize}
}

\changecolours{ScoutGreen}
\section[Referências]{Referências Bibliográficas Recomendadas}

% ------------------------------------------------------------------------------
\begin{frame}{Referências Bibliográficas}
  \itemseps[topsepi=0.5mm,itemsepi=1mm,labelsep=2mm,leftmargini=3.5mm]
  \begin{itemize}
    \scriptsize
    \item \textbf{CHEN, Peter Pin-Shan.} \textit{The Entity-Relationship Model: Toward a Unified View of Data}. ACM Transactions on Database Systems (TODS), New York, v. 1, n. 1, p. 9--36, Mar. 1976.
    \item \textbf{ELMASRI, Ramez; NAVATHE, Shamkant B.} \textit{Sistemas de Banco de Dados}. 7. ed. São Paulo: Pearson Education do Brasil, 2018. (Capítulo 3: Modelagem Conceitual e Capítulo 4: O Modelo EER).
    \item \textbf{HEUSER, Carlos Alberto.} \textit{Projeto de Banco de Dados}. 6. ed. Porto Alegre: Bookman, 2009. (Capítulo 2: O Modelo Entidade-Relacionamento e Capítulo 3: Generalização e Especialização).
    \item \textbf{SILBERSCHATZ, Abraham; KORTH, Henry F.; SUDARSHAN, S.} \textit{Sistema de Banco de Dados}. 7. ed. Rio de Janeiro: Elsevier, 2020. (Capítulo 6: O Modelo E-R Estendido).
    \item \textbf{DATE, C. J.} \textit{Introdução a Sistemas de Bancos de Dados}. 8. ed. Rio de Janeiro: Elsevier / Campus, 2004.
  \end{itemize}
  \vspace{0.15cm}
  \begin{center}
    \tiny\color{gray}
    Instituto Federal de Educação, Ciência e Tecnologia do Ceará -- Campus Tauá\\
    Curso Superior de Tecnologia em Telemática
  \end{center}
\end{frame}


\end{document}
